% =============================================================
%  Module Description — Databases (CS302-F28)
%  Shown on https://wosm.academy as the course page for CS302-F28.
%  Compile with: pdflatex cs302-module-description.tex
%  Write "TODO(owner): ..." where information is still missing;
%  the website highlights it.
% =============================================================
\documentclass[11pt,a4paper]{article}

\usepackage[utf8]{inputenc}
\usepackage[T1]{fontenc}
\usepackage{lmodern}
\usepackage[margin=2.2cm]{geometry}
\usepackage{booktabs}
\usepackage{tabularx}
\usepackage{array}
\usepackage{enumitem}
\usepackage{xcolor}
\usepackage{titlesec}
\usepackage{fancyhdr}
\usepackage[hidelinks]{hyperref}

% ---------- Module details (edit these) ----------------------
% \ModuleCode must equal the course's `code` in src/content/courses/.
\newcommand{\ModuleTitle}{Databases}
\newcommand{\ModuleCode}{CS302-F28}
\newcommand{\Programme}{Foundation programme}
\newcommand{\Level}{Semester 5}
\newcommand{\Semester}{Fall 2028}
\newcommand{\Credits}{7 ECTS}
\newcommand{\ModuleLeader}{Dr. Ali Khudiyev}
\newcommand{\Email}{ali.khudiyev@wosm.academy}
\newcommand{\OfficeHours}{Monday--Friday, 10:00--17:00 (Baku time, UTC+4)}
\newcommand{\Prerequisites}{MAT101-F26 Discrete Structures; CS103-S27 Data Structures and Algorithms I}

% ---------- Styling ------------------------------------------
\definecolor{accent}{RGB}{0,70,130}
\titleformat{\section}{\large\bfseries\color{accent}}{\thesection.}{0.5em}{}
\titlespacing*{\section}{0pt}{1.2em}{0.5em}
\setlist{itemsep=2pt, topsep=4pt}
\setlength{\parindent}{0pt}
\setlength{\parskip}{4pt}

\pagestyle{fancy}
\fancyhf{}
\lhead{\small \ModuleCode\ -- \ModuleTitle}
\rhead{\small Module Description}
\cfoot{\small \thepage}
\renewcommand{\headrulewidth}{0.4pt}

% =============================================================
\begin{document}

\begin{center}
  {\LARGE\bfseries\color{accent} \ModuleTitle}\\[4pt]
  {\large Module Description}
\end{center}
\vspace{6pt}

% ---------- Summary table (not shown on the website: the course page
% shows these details from the macros above) --------------------------
\begin{tabularx}{\textwidth}{>{\bfseries}l X >{\bfseries}l X}
\toprule
Module code   & \ModuleCode    & Credits      & \Credits \\
Programme     & \Programme     & Level        & \Level \\
Semester      & \Semester      & Prerequisites & \Prerequisites \\
Module leader & \ModuleLeader  & Email        & \href{mailto:\Email}{\Email} \\
Office hours  & \multicolumn{3}{l}{\OfficeHours} \\
\bottomrule
\end{tabularx}

% ---------- Sections (shown on the website) ------------------
\section{Module Overview}
Almost every application stores data that must outlive the program, be shared by
many users and stay consistent when things go wrong. This module teaches how to
design, query and build such databases, with the emphasis on relational databases
and SQL.

The first part is about \emph{using} databases: the relational model, relational
algebra, SQL, conceptual design with entity--relationship models, and normalisation.
The second part is about how a database system \emph{works inside}: storage,
indexing, query processing and optimisation, transactions, concurrency control and
recovery. It draws on the data structures from Data Structures and Algorithms I and
the systems ideas from Operating Systems. The module ends with a look at distributed
and non-relational data stores.

\section{Aims}
\begin{itemize}
  \item Teach the relational model and SQL as the standard way to describe and query data.
  \item Develop your ability to design sound database schemas.
  \item Explain how a database system stores, indexes and queries data efficiently.
  \item Explain how transactions keep data consistent under concurrency and failures.
\end{itemize}

\section{Learning Outcomes}
After completing this course, you will be able to:
\begin{enumerate}[label=LO\arabic*.]
  \item Design database schemas with entity--relationship modelling and translate them into relational schemas with keys and integrity constraints.
  \item Express queries in relational algebra and SQL, including joins, aggregation, subqueries and views.
  \item Use functional dependencies to evaluate schemas and normalise them into 3NF and BCNF.
  \item Explain storage and index structures such as B+ trees and hash indexes, and how queries are evaluated and optimised.
  \item Explain transactions, the ACID properties, serialisability, locking-based concurrency control and log-based recovery.
  \item Build an application in Java that uses a relational database through JDBC.
\end{enumerate}

\section{Syllabus}
\begin{tabularx}{\textwidth}{c X}
\toprule
\textbf{Unit} & \textbf{Topic} \\
\midrule
1  & Introduction: data, databases and database management systems \\
2  & The relational model; keys and integrity constraints \\
3  & Relational algebra \\
4  & SQL I: data definition and basic queries \\
5  & SQL II: joins, aggregation, subqueries and NULL values \\
6  & SQL III: views, constraints and triggers; databases from Java with JDBC \\
7  & Conceptual design: entity--relationship modelling \\
8  & From entity--relationship diagrams to relational schemas \\
9  & Functional dependencies and normal forms: 3NF and BCNF \\
10 & Storage: disks, pages, files and records \\
11 & Indexing: B+ trees and hash indexes \\
12 & Query processing: sorting and join algorithms \\
13 & Query optimisation \\
14 & Transactions and concurrency control: serialisability, two-phase locking \\
15 & Recovery: logging and the write-ahead log \\
16 & Beyond relational databases: distributed and NoSQL data stores \\
17 & Course Review \& Exam \\
\bottomrule
\end{tabularx}

\section{Teaching and Learning Methods}
The module is self-paced. All materials are provided on Moodle.
Each unit has a lecture and exercises. Practical work uses a relational database
system to write SQL queries and to build a small database-backed Java application.

\begin{tabularx}{0.6\textwidth}{X r}
\toprule
\textbf{Activity} & \textbf{Hours} \\
\midrule
Lectures            & 34 \\
Labs / Tutorials    & 36 \\
Independent study   & 130 \\
\midrule
\textbf{Total}      & \textbf{200} \\
\bottomrule
\end{tabularx}

\section{Assessment}
Your knowledge will be assessed throughout the course with quizzes, a database
project and the final exam. The project checks whether you can design a schema,
query it in SQL and use it from an application. The exam also covers how database
systems work internally.

\begin{tabularx}{\textwidth}{X c c c}
\toprule
\textbf{Component} & \textbf{Weight} & \textbf{Due} & \textbf{LOs} \\
\midrule
Quizzes & 20\% & Throughout & LO1--LO5 \\
Database project & 40\% & Throughout & LO1, LO2, LO3, LO6 \\
Final exam & 40\% & End of module & LO1--LO5 \\
\midrule
\textbf{Total} & \textbf{100\%} & & \\
\bottomrule
\end{tabularx}

\section{Reading List}
\textbf{Essential}
\begin{itemize}
  \item A.\ Silberschatz, H.\ F.\ Korth, S.\ Sudarshan: \textit{Database System Concepts}, 7th ed., McGraw-Hill, 2019.
\end{itemize}
\textbf{Recommended}
\begin{itemize}
  \item H.\ Garcia-Molina, J.\ D.\ Ullman, J.\ Widom: \textit{Database Systems: The Complete Book}, 2nd ed., Pearson, 2008.
  \item R.\ Ramakrishnan, J.\ Gehrke: \textit{Database Management Systems}, 3rd ed., McGraw-Hill, 2003.
  \item M.\ Kleppmann: \textit{Designing Data-Intensive Applications}, O'Reilly, 2017.
\end{itemize}

\section{Materials and Licence}
Course materials are published under the
\href{https://creativecommons.org/licenses/by-sa/4.0/}{Creative Commons Attribution-ShareAlike 4.0}
licence unless stated otherwise. See \href{https://wosm.academy/licence/}{wosm.academy/licence} for details.

\section{Policies}
\begin{itemize}
  \item \textbf{Attendance:} Attendance is not mandatory. The module is self-paced, so you can study the materials when it suits you.
  \item \textbf{Late submissions:} Work submitted after the deadline is penalised.
  \item \textbf{Academic integrity:} All submitted work must be your own.
        See the \href{https://wosm.academy/code-of-conduct/}{Code of Conduct}.
\end{itemize}

\end{document}
